\documentclass[aspectratio=169, 10pt]{beamer}
\usetheme{Madrid}

% --- Lingua e Pacchetti Base ---
\usepackage[utf8]{inputenc}
\usepackage[italian]{babel}
\usepackage{graphicx}
\usepackage{booktabs}
\usepackage{tcolorbox}
\usepackage{tikz}
\usepackage[colorlinks=true, urlcolor=blue]{hyperref}
\usetikzlibrary{shapes,arrows,positioning}

% --- Personalizzazione dei Colori (Tema KarateFlow) ---
% Colori ispirati al Karate (Rosso cintura, Antracite profondo della divisa/tatami)
\definecolor{karateRed}{HTML}{D90429}
\definecolor{karateDark}{HTML}{2B2D42}
\definecolor{karateLightDark}{HTML}{1F2421}
\definecolor{karateGray}{HTML}{F8F9FA}
\definecolor{karateAccent}{HTML}{8D99AE}

% Configurazione dei colori di Beamer
\setbeamercolor{palette primary}{bg=karateDark, fg=white}
\setbeamercolor{palette secondary}{bg=karateRed, fg=white}
\setbeamercolor{palette tertiary}{bg=karateLightDark, fg=white}
\setbeamercolor{palette quaternary}{bg=karateDark, fg=white}
\setbeamercolor{structure}{fg=karateRed}
\setbeamercolor{titlelike}{parent=palette primary}
\setbeamercolor{block title}{bg=karateDark, fg=white}
\setbeamercolor{block body}{bg=karateGray, fg=black}

% --- Modifiche di Stile ---
\setbeamertemplate{blocks}[rounded][shadow=false]
\setbeamertemplate{navigation symbols}{}
\setbeamertemplate{headline}{}

% Customizzare il piè di pagina (footer) per un look più pulito
\setbeamertemplate{footline}{
  \leavevmode%
  \hbox{%
  \begin{beamercolorbox}[wd=.333333\paperwidth,ht=2.25ex,dp=1ex,center]{author in head/foot}%
    \usebeamerfont{author in head/foot}\insertshortauthor~~(\insertshortinstitute)
  \end{beamercolorbox}%
  \begin{beamercolorbox}[wd=.333333\paperwidth,ht=2.25ex,dp=1ex,center]{title in head/foot}%
    \usebeamerfont{title in head/foot}\insertshorttitle
  \end{beamercolorbox}%
  \begin{beamercolorbox}[wd=.333333\paperwidth,ht=2.25ex,dp=1ex,right]{date in head/foot}%
    \usebeamerfont{date in head/foot}\insertshortdate{}\hspace*{2em}
    \insertframenumber{} / \inserttotalframenumber\hspace*{2ex} 
  \end{beamercolorbox}}%
  \vskip0pt%
}

% --- Dettagli del Titolo ---
\title[KarateFlow]{KarateFlow}
\subtitle{Gestione e Tracciamento delle Performance Atletiche nel Karate}
\author[Simone Ingenito]{Simone Ingenito}
\institute[SPME]{Corso di Software Project Management and Evolution}
\date[Giugno 2026]{23 Giugno 2026}

\titlegraphic{
\includegraphics[width=2.5cm]{images/Karate-flow-logo.png}}

\begin{document}

% Slide 1: Titolo
\begin{frame}[plain]
  \titlepage
\end{frame}

% Slide 2: Contesto
\begin{frame}{1. Il Contesto di Riferimento}
  \begin{columns}
    \begin{column}{0.6\textwidth}
      \textbf{Il Problema delle Associazioni Sportive:}
      \begin{itemize}
        \item Mancanza di strumenti digitali per tracciare la crescita atletica.
        \item Performance registrate su fogli cartacei disordinati.
        \item Difficoltà nel confrontare storicamente le performance degli atleti.
      \end{itemize}
      \vspace{0.3cm}
      \textbf{L'Esigenza Specifica:}
      \begin{itemize}
        \item Fornire al preparatore e ai coach uno storico centralizzato.
        \item Tracciare periodicamente parametri fisici e tecnici reali.
      \end{itemize}
    \end{column}
    \begin{column}{0.35\textwidth}
      \begin{tcolorbox}[colback=karateRed!10, colframe=karateRed, title=Inquadramento]
        Digitalizzazione e automazione dei flussi di lavoro in un'associazione sportiva locale per migliorare la gestione quotidiana degli atleti sul tatami.
      \end{tcolorbox}
    \end{column}
  \end{columns}
\end{frame}

% Slide 3: Motivazioni
\begin{frame}{1. Le Motivazioni del Progetto}
  \begin{columns}
    \begin{column}{0.6\textwidth}
      \textbf{La Soluzione: KarateFlow}
      \begin{itemize}
        \item Web App dedicata per inserimento rapido dei test atletici periodici.
        \item Centralizzazione dell'anagrafica atleti e registrazione dei test.
        \item Confronto e misurazione delle prestazioni.
      \end{itemize}
    \end{column}
    \begin{column}{0.35\textwidth}
      \begin{tcolorbox}[colback=karateRed!10, colframe=karateRed, title=Obiettivo Primario]
        Creare una piattaforma su misura per il mio team di karate, mantenendo i costi bassi.
      \end{tcolorbox}
    \end{column}
  \end{columns}
\end{frame}

    % Slide: Studente Atleta                                                                                                           
    \begin{frame}{1. Note sull'Autore: Studente e Atleta}                                                                            
      \begin{columns}                                                                                                                     
        \begin{column}{0.5\textwidth}
          \textbf{La Doppia Carriera:}
          \begin{itemize}
            \item Coniugare lo studio accademico in Informatica con la pratica agonistica del Karate.
            \item Membro del programma \href{https://atletafedericiano.unina.it/}{Atleta Federiciano}  presso l'Università degli Studi di Napoli Federico II.
            \item Road to \textbf{EUSA Championships}, Tirana Settembre 2026.     
          \end{itemize}
          \vspace{0.3cm}
          \textbf{Il Collegamento con KarateFlow:}
          \begin{itemize}
            \item Il software nasce da un'esigenza personale e del mio team di karate.
          \end{itemize}
        \end{column}
        \begin{column}{0.45\textwidth}
          \centering
          \begin{minipage}{0.48\linewidth}
            \centering
            \includegraphics[height=3.8cm, width=\linewidth, keepaspectratio]{images/CNU-2026.jpg}\\
            \vspace{0.1cm}
            \tiny{Campionati Nazionali Universitari, Maggio 2026}
          \end{minipage}
          \hfill
          \begin{minipage}{0.48\linewidth}
            \centering
            
\includegraphics[height=3.8cm, width=\linewidth, keepaspectratio]{images/atleta-federiciano.jpg}\\
            \vspace{0.1cm}
            \tiny{Programma Atleta Federiciano}
          \end{minipage}
        \end{column}
      \end{columns}
    \end{frame}

% Slide 4: Perché SPME?
\begin{frame}{2. Perché è stato scelto come progetto SPME?}
  Il corso di \emph{Software Project Management and Evolution} richiede di unire la gestione organizzativa allo sviluppo software moderno.
  \vspace{0.4cm}
  \begin{columns}
    \begin{column}{0.31\textwidth}
      \begin{tcolorbox}[colback=karateDark!5, colframe=karateDark, title=1. Metodologie DevOps, height=4.5cm]
        Sperimentare lo sviluppo software con metodologie DevOps.
      \end{tcolorbox}
    \end{column}
    \begin{column}{0.31\textwidth}
      \begin{tcolorbox}[colback=karateDark!5, colframe=karateDark, title=2. Project Management, height=4.5cm]
        Fare pratica come Project Manager nel contesto dello sviluppo software.
      \end{tcolorbox}
    \end{column}
    \begin{column}{0.31\textwidth}
      \begin{tcolorbox}[colback=karateDark!5, colframe=karateDark, title=3. Traceability, height=4.5cm]
        Fare pratica della gestione del progetto software tramite Ticketing, garantendo il collegamento bidirezionale tra i requisiti e il codice. 
      \end{tcolorbox}
    \end{column}
  \end{columns}
\end{frame}

% Slide 5: Il Mio Team: IA come Co-pilota
\begin{frame}{3. Il Mio Team: Sviluppo Assistito da IA}
  Lavorando come \textbf{unico sviluppatore}, ho massimizzato l'efficienza utilizzando \textbf{Gemini} come agente per la scrittura rapida del codice, interpretando diversi ruoli:
  \vspace{0.2cm}
  \begin{columns}
    \begin{column}{0.48\textwidth}
      \begin{block}{I Ruoli Simulati dall'Agente}
        \begin{itemize}
          \item \textbf{Backend Engineer:} Scrittura del codice Spring Boot e modellazione del DB.
          \item \textbf{Frontend Engineer:} Sviluppo dei componenti Angular responsive.
          \item \textbf{DevOps Specialist:} Supporto alla definizione delle pipeline Jenkins, setup di Terraform e Kubernetes.
        \end{itemize}
      \end{block}
    \end{column}
    \begin{column}{0.48\textwidth}
      \begin{block}{Benefici \& Metodologia}
        \begin{itemize}
          \item \textbf{Velocità:} Sviluppo di scheletri strutturali in pochi minuti.
          \item \textbf{Qualità:} Riduzione dei bug logici grazie a prompt orientati ai test di unità.
          \item \textbf{Rigore:} Prima di ogni modifica remota, il codice è stato sottoposto ad analisi assicurando standard di qualità eccellenti.
        \end{itemize}
      \end{block}
    \end{column}
  \end{columns}
  \begin{itemize}
  \vspace{0.2cm}
      \item Nota: Ho riscattato il piano pro ad un prezzo molto conveniente (2 euro al mese per tre mesi) e cosi ho potuto simulare lo sviluppo software di un team ben strutturato, concentrandomi sul ruolo di PM.
  \end{itemize}
\end{frame}

% Slide 6: Architettura a blocchi
\begin{frame}{4. Architettura e Stack Tecnologico}
  \begin{center}
    \begin{tikzpicture}[
      node distance=1.2cm and 1cm,
      box/.style={draw, rectangle, rounded corners, minimum width=2.5cm, minimum height=0.8cm, align=center, fill=karateGray, font=\footnotesize},
      accentbox/.style={draw=karateRed, rectangle, rounded corners, minimum width=2.5cm, minimum height=0.8cm, align=center, fill=karateRed!10, font=\footnotesize},
      arrow/.style={-latex, thick}
    ]
      % Nodes
      \node[box] (fe) {\textbf{Frontend (Angular)}\\SPA Statica (Nginx)};
      \node[box, right=of fe] (be) {\textbf{Backend (Spring Boot)}\\REST API};
      \node[box, right=of be] (db) {\textbf{Database (MongoDB)}\\Dati Flessibili NoSQL};
      \node[accentbox, below=of be] (k3s) {\textbf{Cluster K3s}\\Raspberry Pi 5 Host (Staging \& Prod)};
      \node[box, left=of k3s] (infra) {\textbf{IaC (Terraform)}\\Definizione Limiti/Namespace};
      \node[box, right=of k3s] (tunnel) {\textbf{Cloudflare Tunnel}\\Connettività HTTPS Sicura};

      % Arrows
      \draw[arrow] (fe) -- (be);
      \draw[arrow] (be) -- (db);
      \draw[arrow] (infra) -- (k3s);
      \draw[arrow] (k3s) -- (tunnel);
      \draw[arrow] (fe) .. controls +(down:1) and +(left:0.5) .. (k3s);
      \draw[arrow] (be) -- (k3s);
    \end{tikzpicture}
  \end{center}
  \vspace{0.2cm}
  \centering
  \scriptsize{L'architettura multi-repository si basa su servizi isolati ed integrati in un cluster kubernetes.}
\end{frame}

% Slide 7: Hosting e Vincoli Fisici
\begin{frame}{4. Hosting e Vincoli Fisici}
  \begin{columns}
    \begin{column}{0.6\textwidth}
      \textbf{L'Hosting su Raspberry Pi 5:}
      \begin{itemize}
        \item Consumi ridotti al minimo per azzerare i costi dell'hosting in cloud.
        \item Ideale per una piccola realtà sportiva locale.
        \item \alert{I Vincoli Hardware:} Risorse RAM e CPU estremamente limitate.
      \end{itemize}
      \vspace{0.3cm}
      \textbf{Trade-off e Ottimizzazioni di Stack:}
      \begin{itemize}
        \item Compilazione statica in Angular (servita via Nginx) per minimizzare la memoria RAM occupata in esecuzione.
        \item Uso di container leggeri privi di componenti non necessari.
        \item Resource limits configurati rigidamente per evitare memory leak distruttivi per l'host.
      \end{itemize}
    \end{column}
    \begin{column}{0.35\textwidth}
      \begin{tcolorbox}[colback=karateDark!5, colframe=karateDark, title=Specifiche Host]
        \textbf{Hardware:} Raspberry Pi 5\\
        \textbf{RAM:} 8 GB\\
        \textbf{Storage:} SSD USB 3.0\\
        \textbf{OS:} Raspberry Pi OS (Debian 12 Bookworm)\\
      \end{tcolorbox}
    \end{column}
  \end{columns}
\end{frame}

% Slide 8: IaC e Orchestrazione (Terraform e Kubernetes)
\begin{frame}{5. IaC e Orchestrazione (Terraform + K3s)}
  \begin{columns}
    \begin{column}{0.52\textwidth}
      \textbf{Virtualizzazione Leggera con K3s:}
      \begin{itemize}
        \item Distribuzione Kubernetes ottimizzata per IoT ed Edge Computing.
        \item Consumo di memoria ridotto per il Control Plane.
      \end{itemize}
      \vspace{0.3cm}
      \textbf{Il Ruolo Centrale di Terraform (IaC):}
      \begin{itemize}
        \item Configurazione programmatica dei namespace (\texttt{stg} e \texttt{prod}).
        \item Definizione di \alert{Resource Quotas} rigide (limiti CPU/RAM) per isolare gli ambienti e prevenire il crash totale del server domestico.
        \item \textbf{Disaster Recovery:} Ricreazione totale del cluster in poco tempo in caso di guasti.
      \end{itemize}
    \end{column}
    \begin{column}{0.44\textwidth}
      \begin{tcolorbox}[colback=karateDark!5, colframe=karateDark, title=Resource Isolation Framework, left=1mm, right=1mm, top=1mm, bottom=1mm]
        \centering
        \scriptsize
        \resizebox{\linewidth}{!}{
          \begin{tabular}{lcc}
            \toprule
            \textbf{Namespace} & \textbf{CPU Limit} & \textbf{RAM Limit} \\
            \midrule
            \texttt{karateflow-stg} & 1.5 Cores & 2.0 GB \\
            \texttt{karateflow-prod} & 2.5 Cores & 3.5 GB \\
            \bottomrule
          \end{tabular}
        }
      \end{tcolorbox}
    \end{column}
  \end{columns}
\end{frame}

% Slide 9: Navigazione dei Repository
\begin{frame}{6. Navigazione del Repository di Documentazione}
    
    Ho creato un'organizzazione Github che racchiude tutto: \url{https://github.com/KarateFlow}.

  \vspace{0.2cm}
  Il repository \textbf{\texttt{karateflow-docs}} funge da cabina di regia del progetto e archivio di tutte le scelte fatte e dei progressi raggiunti step-by-step:
  \vspace{0.4cm}

      \begin{itemize}
        \item \textbf{\texttt{requisiti-architettura.md}:} Contiene i vincoli del Raspberry Pi 5 e la roadmap evolutiva di alto livello.
        \item \textbf{\texttt{branching-strategy.md}:} Definisce lo sviluppo basato su Trunk-Based Development.
        \item \textbf{\texttt{guida-tracciabilità.md}:} In cui spiego l'approccio utilizzato per la gestione della board e dei Ticket.
        \item \textbf{\texttt{wbs.md}:} Traccia l'elenco gerarchico dei Work Packages, che sono tradotti in EPICHE.
        \item \textbf{Cartella \texttt{diario-di-bordo/}:} File markdown incrementali che descrivono l'andamento del progetto e le decisioni prese in ogni singola fase.
      \end{itemize}

\end{frame}

% Slide 10: Struttura dei Repository BE e FE
\begin{frame}{7. Uno Sguardo alle Repository di Backend e Frontend}
  La struttura delle repository è stata pensata per massimizzare la leggibilità e ottimizzare il contesto dei modelli di IA. 
  \vspace{0.3cm}
  \begin{columns}
    \begin{column}{0.48\textwidth}
      \begin{block}{Backend: Spring Boot (\texttt{karateflow-be})}
        \begin{itemize}
          \item \textbf{Architettura Esagonale (Ports \& Adapters):}
          \item Core Domain (\emph{model} e \emph{port}) slegato da librerie esterne.
          \item Logica applicativa orchestrata da \emph{usecase}.
          \item Persistenza (MongoDB Adapter) e Controller Web come \emph{adattatori} esterni intercambiabili.
          \item Alta testabilità con test unitari e d'integrazione indipendenti.
        \end{itemize}
      \end{block}
    \end{column}
    \begin{column}{0.48\textwidth}
      \begin{block}{Frontend: Angular (\texttt{karateflow-fe})}
        \begin{itemize}
          \item Single Page Application organizzata in moduli feature.
          \item Utilizzo di componenti riutilizzabili e puliti, privi di logica di business.
          \item Separazione netta tra UI e Business Logic. Posso cambiare totalmente interfaccia senza "rompere" l'applicazione.
          \item Compilazione statica e servita tramite Nginx per minimizzare il consumo di RAM del server.
        \end{itemize}
      \end{block}
    \end{column}
  \end{columns}
\end{frame}

% Slide 11: Repository di Infra e Automazione CI/CD
\begin{frame}{8. Repository di Infra e Pipeline CI/CD}
  \begin{columns}
    \begin{column}{0.5\textwidth}
      \textbf{La Pipeline DevOps con Jenkins:}
      \begin{itemize}
        \item Jenkins è installato sulla \alert{macchina di sviluppo locale} (non sul Raspberry Pi 5) per evitare overhead di CPU e RAM.
        \item Trigger di build configurati tramite polling Git o esecuzione manuale.
        \item \textbf{Dockerizzazione Immutabile:} Dockerfile multi-stage che compilano il sorgente ed esportano immagini JRE/Nginx minime e sicure.
        \item \textbf{Integrazione con GitHub Status Checks:} Jenkins "avvisa" GitHub dell'esito dei test, inibendo il merge delle Pull Request in caso di build fallite.
      \end{itemize}
    \end{column}
    \begin{column}{0.48\textwidth}
      \begin{block}{Pipeline Status}
        \centering
        \includegraphics[width=0.9\textwidth]{images/Jenkins-check.png}\\
        \vspace{0.1cm}
        \scriptsize{Pipeline Jenkins per la validazione statica (PMD/ESLint), la compilazione dei test e il build delle immagini Docker.}
      \end{block}
    \end{column}
  \end{columns}
\end{frame}

% Slide 12: Tracciabilità dei Requisiti
\begin{frame}{9. Tracciabilità dei Requisiti}
  \begin{columns}
    \begin{column}{0.5\textwidth}
      \textbf{Flusso di Tracciabilità Bidirezionale:}
      \begin{enumerate}
        \item I requisiti vengono tradotti in EPIC/STORY e aggiunte alla Kanban Board. 
        \item Le Stories contengono scenari scritti in formato \textbf{BDD (Given-When-Then)}.
        \item Il branch di sviluppo viene creato e collegato alla issue presente sulla board (Task o Bug).
        \item La Pull Request verso \texttt{main} chiude l'issue tramite Closing Keyword (\texttt{Closes \#12}) sul server.
        \item Nessun frammento di codice viene inserito senza una motivazione di business documentata.
      \end{enumerate}

    \end{column}
    \begin{column}{0.34\textwidth}
      \begin{block}{Esempio di PR}
        \centering
        \includegraphics[width=0.7\textwidth]{images/template-PR.png}\\
        \vspace{0.1cm}
        \scriptsize{Collegamento bidirezionale tra Pull Request e Issue.}
      \end{block}
    \end{column}
  \end{columns}
\end{frame}

% Slide 13: GitHub Projects vs Jira
\begin{frame}{9. Governance e Tracciamento: GitHub Projects vs Jira}
 
  \vspace{0.3cm}
  \begin{columns}
    \begin{column}{0.95\textwidth}
      \begin{tcolorbox}[colback=karateGray, colframe=karateDark, title=Confronto Strumenti di Project Management]
        \centering
        \small
        \resizebox{\linewidth}{!}{
          \begin{tabular}{lll}
            \toprule
            \textbf{Caratteristica} & \textbf{GitHub Projects} & \textbf{Jira Software} \\
            \midrule
            \textbf{Integrazione Git} & Nativa e immediata & Tramite plugin esterni \\
            \textbf{Overhead Amministrativo} & Minimo & Elevato (config. workflow) \\
            \textbf{Tracciabilità Ticket-Codice} & Automatica (issue $\leftrightarrow$ branch $\leftrightarrow$ PR) & Richiede inserimento manuale ID \\
            \textbf{Adatto per} & Sviluppatore singolo o piccoli team & Grandi aziende e team complessi \\
            \bottomrule
          \end{tabular}
        }
      \end{tcolorbox}
      \vspace{0.2cm}
      \textbf{Perché GitHub Projects?}
      \begin{itemize}
        \item Per KarateFlow, Jira avrebbe introdotto un overhead amministrativo inutile.
        \item L'integrazione nativa permette di collegare facilmente ticket e codice.
      \end{itemize}
    \end{column}
  \end{columns}
\end{frame}

% Slide 14: Prima Fase di Sviluppo e Gestione Task
\begin{frame}{10. Metodologia e Prima Fase di Sviluppo}
  \begin{columns}
    \begin{column}{0.5\textwidth}
      \textbf{Punti di Forza Metodologici:}
      \begin{itemize}
        \item \textbf{Trunk-Based Development:} Rilasci continui sul branch \texttt{main} (allineato all'ambiente di Staging).
        \item \textbf{Conventional Commits:} Standardizzazione dei messaggi Git (\texttt{feat}, \texttt{chore}, \texttt{docs}, \texttt{devops}).
        \item \textbf{Rimozione delle Stime Orarie:} Per non introdurre un overhead amministrativo inutile per un singolo sviluppatore, ci si è concentrati sull'atomicità dei task (tutti dimensionati tra le 2 e le 4 ore di lavoro ideale).
      \end{itemize}
    \end{column}
    \begin{column}{0.48\textwidth}
      \begin{block}{Stato della Kanban Board (inizio fase 4)}
        \centering
        \includegraphics[width=0.9\textwidth]{images/board-inizio-fase-4.png}\\
        \vspace{0.1cm}
        \scriptsize{La Kanban Board di GitHub all'inizio della Fase 4 prima dello sviluppo dei task di automazione.}
      \end{block}
    \end{column}
  \end{columns}
\end{frame}

% Slide 15: Retrosettiva Fase 3: Difficoltà e Soluzioni
\begin{frame}{10. Retrosettiva Fase 3: Difficoltà e Soluzioni}
  Durante la fase di sviluppo core (MVP) sono emerse alcune inefficienze operative:
  \vspace{0.3cm}
  \begin{columns}
    \begin{column}{0.48\textwidth}
      \begin{block}{Le Difficoltà Riscontrate}
        \begin{itemize}
          \item \textbf{Story Mismatch / Scope Creep:} Alcune User Story si sono allargate a dismisura e hanno superato il loro scope, a causa di pipeline CI/CD configurate male o incomplete che ritardavano il feedback.
          \item \textbf{Inconsistenza dei Task:} I nomi dei task e dei commit su Git non erano uniformi perché gestiti manualmente a mente sul portale.
          \item \textbf{Inconsistenza delle PR:} Inizialmente scrivevo ogni PR manualmente e la struttura delle varie PR era inconsistente.
        \end{itemize}
      \end{block}
    \end{column}
    \begin{column}{0.48\textwidth}
      \begin{block}{Le Soluzioni Introdotte}
        \begin{itemize}
          \item \textbf{Issue Templates:} Standardizzazione rigida della creazione di Task, Bug e Requisiti in formato BDD.
          \item \textbf{Pull Request Templates:} Standardizzazione del formato di una PR con descrizione delle modifiche, test eseguiti, possibile immagine se prevista e sezione di tracciabilità con chiusura automatica del TASK associato.
          \item \textbf{Integrazione Branch-Issue:} I branch di sviluppo vengono eliminati una volta effettuata la PR.
        \end{itemize}
      \end{block}
    \end{column}
  \end{columns}
\end{frame}

% Slide 16: Jump-off alla fine della fase 4
\begin{frame}{11. Jump-off alla Fine della Fase 4 (Milestone Staging)}
  \begin{columns}
    \begin{column}{0.5\textwidth}
      \textbf{Stato dell'Infrastruttura Raggiunto:}
      \begin{itemize}
        \item Provisioning del cluster K3s su Raspberry Pi 5 completato tramite Terraform.
        \item Isolamento completo degli ambienti tramite namespace e quote.
        \item Pipeline CI/CD di Jenkins operative ed automatizzate al 100\% post-merge.
      \end{itemize}
      \vspace{0.3cm}
      \textbf{Accesso Remoto e Sicurezza:}
      \begin{itemize}
        \item Utilizzo di un \alert{Cloudflare Tunnel} (\texttt{cloudflared}) nel cluster: comunicazione crittografata senza fare Port Forwarding sul router di casa.
        \item Registrazione e collegamento del dominio di proprietà: \texttt{karate-flow.com}.
      \end{itemize}
    \end{column}
    \begin{column}{0.48\textwidth}
      \begin{tcolorbox}[colback=karateRed!10, colframe=karateRed, title=Staging Online]
        L'applicazione in ambiente di staging è ora accessibile pubblicamente tramite connessione sicura HTTPS al link:
        \vspace{0.3cm}
        \centering
        \textbf{\url{https://staging.karate-flow.com}}
      \end{tcolorbox}
    \end{column}
  \end{columns}
\end{frame}

% Slide 17: La Fase 5: Nuove EPIC
\begin{frame}{12. La Fase 5: Le Nuove EPIC (Sviluppi in Staging)}
  Completate le fondamenta infrastrutturali, il progetto prevede l'implementazione di nuove funzionalità di governance, sicurezza, monitoraggio e miglioramento dell'interfaccia utente.
  \vspace{0.3cm}
  \begin{columns}
    \begin{column}{0.48\textwidth}
      \begin{block}{Autenticazione \& Sicurezza (Keycloak)}
        \begin{itemize}
          \item Integrazione di Keycloak nel cluster K3s per gestire l'autenticazione centralizzata.
          \item Definizione di ruoli RBAC (Role-Based Access Control) per Coach, Atleta e Admin.
        \end{itemize}
      \end{block}
      \begin{block}{Qualità del Codice (SonarQube)}
        \begin{itemize}
          \item Integrazione di SonarQube nelle pipeline di Jenkins per il controllo statico della qualità e delle vulnerabilità di sicurezza.
        \end{itemize}
      \end{block}
    \end{column}
    \begin{column}{0.48\textwidth}
      \begin{block}{Monitoraggio \& Log (Stack PLG)}
        \begin{itemize}
          \item Installazione di Promtail, Loki e Grafana (PLG Stack) su K3s.
          \item Centralizzazione dei log dei pod di backend, frontend e database per semplificare il debug in produzione.
        \end{itemize}
      \end{block}
      \begin{block}{Miglioramenti UI/UX}
        \begin{itemize}
          \item Refactoring dell'interfaccia Angular per una dashboard atletica avanzata e visualizzazione di grafici storici.
        \end{itemize}
      \end{block}
    \end{column}
  \end{columns}
\end{frame}

% Slide 18: La Fase 6: Messa in Produzione
\begin{frame}{13. La Fase 6: Messa in Produzione}
  L'ultimo step prevede il rilascio stabile in ambiente di produzione e l'utilizzo del software da parte degli utenti finali.  
  \vspace{0.4cm}
  \begin{columns}
    \begin{column}{0.5\textwidth}
      \textbf{Strategia di Release:}
      \begin{itemize}
        \item Deployment nel namespace \texttt{karateflow-prod} del cluster K3s locale.
        \item Promozione del codice tramite Tag Git (es. \texttt{v1.0.0}) anziché push continuo.
        \item Isolamento hardware completo gestito tramite le quote Terraform.
      \end{itemize}
    \end{column}
    \begin{column}{0.48\textwidth}
      \begin{block}{Attività Chiave di Produzione}
        \begin{itemize}
          \item \textbf{Load Testing:} Validare il comportamento del Raspberry Pi 5 sotto stress di rete.
          \item \textbf{Dati Reali:} Importazione dell'anagrafica e dei test del team.
          \item \textbf{SSL e DNS:} Routing del dominio principale (\url{https://karate-flow.com}) verso l'ambiente di produzione.
        \end{itemize}
      \end{block}
    \end{column}
  \end{columns}
\end{frame}

% Slide 19: Sviluppi Futuri
\begin{frame}{14. Sviluppi Futuri}
 La piattaforma KarateFlow è predisposta per un'espansione e miglioramento continuo:
  \vspace{0.4cm}
  \begin{columns}
    \begin{column}{0.48\textwidth}
      \begin{tcolorbox}[colback=karateDark!5, colframe=karateDark, title=1. Migrazione Cloud Semplificata, height=4.5cm]
        Grazie all'infrastruttura interamente dichiarata in Terraform e manifesti Kubernetes, è possibile migrare l'intera piattaforma su un cloud provider (AWS, GCP o altri) con facilità senza riscrivere configurazioni hardware.
      \end{tcolorbox}
    \end{column}
    
         \begin{column}{0.48\textwidth}
          \begin{tcolorbox}[colback=karateDark!5, colframe=karateDark, title=2. Monetizzazione Futura (SaaS), height=4.5cm]
            \begin{itemize}
              \item Modello \textbf{SaaS} da vendere a potenziali clienti nell'ambito sportivo. 
              \item Piani di Abbonamento con funzionalità differenziate.
              \item Payment Gateway con gestione fatturazioni ricorrenti e pagamenti automatici.
            \end{itemize}
          \end{tcolorbox}
        \end{column}
    
  \end{columns}
\end{frame}

% Slide 20: Conclusione
\begin{frame}{15. Conclusione}
  \begin{columns}
    \begin{column}{0.6\textwidth}
      \textbf{Risultati Raggiunti:}
      \begin{itemize}
        \item Realizzazione di una piattaforma web completa ed operante in staging.
        \item Applicazione delle metodologie DevOps e di Project Management.
        \item Tracciabilità bidirezionale implementata (con qualche errore iniziale).
      \end{itemize}
      \vspace{0.3cm}
      \textbf{Cosa ho imparato:}
      \begin{itemize}
        \item Metodologia DevOps, rilascio continuo di nuove funzionalità software. 
        \item Utilizzo di Infrastructure as Code (Terraform) e Orchestrazione dei container (Kubernetes).
        \item Utilizzo di strumenti AI per lo sviluppo di codice.
      \end{itemize}
    \end{column}
    \begin{column}{0.35\textwidth}
      \begin{tcolorbox}[colback=karateDark, colframe=karateDark, title=\color{white}Grazie per l'Attenzione]
        \color{white}
        \centering
        \vspace{0.5cm}
        \textbf{Domande?}
        \vspace{0.8cm}
      \end{tcolorbox}
    \end{column}
  \end{columns}
\end{frame}

\end{document}
